\documentclass[]{../../../../tex/classes/styledArticle}
\usepackage{../../../../tex/packages/hebrewSupport}


\author{שחר פרץ}
\title{\textit{היסטוריה 49}}
\begin{document}
	\maketitle
	
	
	בחינה מתחילה ברבע ל־8 ונערכת במתכונת של בגרות. בודקים לאיזה חדר שמשובצים אליו והולכים אל החדר המתאים. יש מקסימום 18 ילדים בחדר, בחלקם פחות. כל ילד על שולחן נפרד. מקבלים מחברות בחינה ומתייחסים אליהן כמו מחברות בגרות – לא תולשים, כותבים שורה כן ושורה לא (אלו מחברות משובצות) אחרת אי אפשר לקרוא מה כותבים. אם לוקחים מחברות נוספת, תכתבו שיש מחברת נוספת על המחברת הראשונה, ועל השנייה לכתוב שיש מחברת ראשונה. אורך 3 שעות + תוספת זמן (בחדר יעודי). 
	
	זוכרים שבפרק 4 היו 3 שאלות? יש 4. המורה יכול לבחור האם לומדים יום הכיפורים או ששת הימים. יש מקרים בהם השאלה על מלחמות ישראל היא שאלה אחת, ובה בוחרים מלחמה ועונים בהתאם למלחמה עליהם למדתם. בהינתן המיקוד הנוכחי כנראה תהיה שאלה על מלחמת יום כיפור, על אף שכבר בעבר היו מקרים בהם לא הייתה שאלה על יום כיפור. 
	
	שעת סיפור: המורה של הבת של שרית המרה והחליטה ללמד רק יום כיפור. יומיים לפני הבגרות שרית הושיבה אותה ועוד 4 מחברותיה ולמדה אותן תוך שעה את דה־קולוניזציה. לא הייתה שאלה על יום כיפור בבגרות וזה מה שהציל להן את הבגרות. 
	
	בשיעור הזה נחזור על יום כיפור ונלמד את המלחמה יותר לעומק. מלחמת יום הכיפורים מתחלקת לשני חלקים, חלק א' וחלק ב', והם מחולקים בתורם לתקופה 1, 2, 3, 4. 
	
	בתקופה 2 בתוך חלק א' יש לנו את תוכנית ד', תקופה 4 בתוך חלק ב' נפתחת בקרבות עשרת הימים ונמשכת אל תוך והמבצעים הגדולים – יואב, חירם, עובדה ועוד משהו אני לא מצליח לקרוא מהלוח. 
	
	\subsection*{תקופות 1 ו־2}
	
	המציאות שנוצרה אחרי כ''ט בנובמבר יוצרת כמה דברים – גושי התיישבות, בעיות בדרכים, יישובים מבודדים, ירושלים נצורה, ערים מעורבות בעייתיות, וכו'. אם מדברים על הקשיים בראשית המלחמה, הם נובעים באופן ישיר מתוכנית החלוקה עצמה. מה שמאוד מאפיין את התקופה הזו הן השיירות, שרק חלקן מצליחות לעבור. בינואר יש לנו את הסיפור של שיירת הל''ה, אירוע מדגים שמספר על הקשיים של התקופה הזו. דיברנו עליה בשיעור הקודם של מלחמת יום העצמאות. 
	
	כל המלחמה, על כל חלקיה, היא מלחמת קיום – להיות או לחדול. מדברים על המלחמה הזו כ\textit{''המלחמה המעצבת``} – היא זו שעצבה את גבולות המדינה, חלק גדול מהערכים בצה''ל, חלק לא קטן במערכת היחסים בין השלטון לצב''א, וכו'. 
	
	במהלך התקופה הראשונה, ארה''ב הציעה כל מני חלופות, כלומר נסוגה מתמיכתה בתוכנית החלוקה. ההצעה: משטר נאמנות של האו''ם, עד ''לא יודעים``. שי אי־הסכמות גם בתוך המשטר האמריקאי בנושא הזה. יש מפגש ל וייצמן עם טרומן, עם כל מני הבטחות מתחת לשולחן. הנהגת היישוב לא יכולה לקבל את העמדה הרשמית. 
	
	בתקופה בה תוכנית ד' יצאה אל הפועל כבר היה ברור שהבריטים יוצאים ב־14 למאי, והיה ברור שחייבים לשנות את המצב הצבאי, ולא רק זה – גם את המבנה הצבאי. ''אין אף מטומטם שחושב שפה (מצביעה על יציאת הבריטים) לא תפרוץ מלחמה עם מדינות ערב``, ובמבנה הנוכחי של המחתרות זה לא עומד לעבוד. תוכנית ד' מתחילה את התקופה השנייה בחלק א'. 
	
	בתוכנית ד' גם מתחילה התארגנות של לבנות את הצבא. מתחילים לבנות יחידות, פיקוד של אזורים, וכו'. חלק גדול מהקמת צה''ל קורה בחודש של תוכנית ד' / תקופה 2. המטרות: 
	\begin{multicols}{2}
		\begin{itemize}
			\item שליטה על ערים מעורבות
			\item רצף טריטוריאלי
			\item שליטה על יישובים מבודדים
			\item לפתוח את הדרך לירושלים
		\end{itemize}
	\end{multicols}
	ובתוך זה, כל מה שקשור ליישובים הערבים – אם הם מפריעים לרצף הטריטוריאלי, צריך לכבוש אותם. בנקודת הזמן הזו מתחילה הבריחה וכן בעיית הפליטים. גירוש, יש, פחות בתקופה של תוכנית ד'. רוב הפלסטינים בורחים ולא מגורשים (באופן כללי, לא רק בתוכנית ד'). 
	
	למה קוראים לה ד'? כי לפני כן היה א', ב', וג'. 
	
	לגבי אירוע שמאפיין את התקופה – יש לנו את מבצע נחשון. ירושלים אמורה להלחם בירדן, צבא חזק. לאנשים בירושלים אין אוכל ונשק. בן גוריון (''ממש לא בן־אדם דתי, אין להבנתו שום היבט דתי``) מבין את חשיבותה גם בהבט המורלי, ומחליט לפרוץ אותה ויהי מה. בן גוריון נותן את הפקודה לרכז 1500 חיילים כדי לפרוץ את הדרך לירושלים, על חשבון אזורים אחרים (ולמקומות אחרים לא הייתה הגנה). הצליחו לפרוץ את הדרך לירושלים ולהביא לשם שיירות גדולות. אחרי זמן קצר, הערבים כובשים מחדש את האיזורים הגבוהים, והדרך נחסמת. זמן קצר אחרי הכרזת העצמאות הצליחו לפתוח את דרך בורמה, דרך עוקפת שדרכה יצרו קשר לירושלים. 
	
	אם מדברים על המאפיינים של תוכנית ד' – ריכוז כוחות, יוזמה התקפית, בנייה של כוח לחימה וכו', רואים זאת במבצע נחשון. 
	
	\subsection*{תקופות 3 ו־4}
	הבסיס לצה''ל וההגנה, היחידות הלוחמות הסדירות, מתחילות מהפלמ''ח, שבן גוריון מאוד מהר מפרק (לא נכנס כרגע ללמה, יש שם גם עניינים אידיאולוגים). ביוני 1948 יש את הפקודה להקמת צה''ל. אותה השבועה מאז נגררת עד היום. בן גוריון מקבל ערך שנקרא ''ערך הממלכתיות``. בתהליך הזה של בניית צה''ל יש תהליך של פירוק המחתרות. בדרך יש כל מני מהמורות, אחת מהן אלטלנה.
	
	אלטלנה – ספינה שמביא האצ''ל, עוד לפני עם נשק שנרכש לפני שהוקמה המדינה. יש מעפילים את הספינה הזו. בשם הממלכתיות, צבא אחד, מדינה אחת, האצ''ל נדרש לתת את הנשק לצה''ל. לאצ''ל יש כוחות בירושלים, והיא לא חלק מהמדינה, ולכן הכוחות לא התפרקו שם – אין שם צה''ל, האצ''ל נשאר שם מחתרת. הם דורשים לתת אחוז מסוים מהנשק לאצ''ל בירושלים. בן־גוריון טוען שאין להם זכות לדרוש דבר כזה. נערך שם מאבק בסיומו צה''ל מפציץ את הספינה לאחר מאבק מזוין עם האצ''ל, נהרגו איזה 15 איש ונהרסו עשרות נשקים. 
	
	סוריה, לבנון, עיראק (כוח עזר)'' ירדן ומצרים – פולשים לארץ מתוך תפיסה שעושים בליצקריג/ונצואלה לארץ, בשנייה שהבריטים יצאו. זה היה ברור שזה יקרה – עוד באו''ם מדינות ערב אמרו שברגע שתוכרז כאן מדינה יפלשו לארץ. ישראל נמצאת במגננה. דיברנו בשיעורים קדומים על הטיעונים בעד ונגד ההכרזה. 
	
	אם יהיה לכם זמן, תלמדו בבית את הנושא של פירוק מטה הפלמ''ח. יש ויכוח בעד ונגד פירוק מטה הפלמ''ח. מדינה ריבונית לא יכולה להרשות לעצמה שיהיו לה כוחות צבאיים עצמאיים, בד''כ גם יש להם אג'נדה פוליטית. בן־גוריון בשם הממלכתיות, אומר שיש מדינה אחת, ומחליט לפרק לא רק את המחתרות הפולשות, אלא גם את הפלמ''ח. יש כאן גם עניין אידיאולוגי – הפלמ''ח זה לא מפא''י, ויש הטוענים שזו הסיבה שבן־גוריון החליט לפרק את מטה הפלמ''ח. יש בעייתיות מעשית – באמצע מלחמה, צבאות מדינות ערב פולשים, ובן גוריון מפרק את היחידות שהיוו את הבסיס הלוחם הרציני. יש מאבק בין בן־גוריון לבין מטה הפלמ''ח. בעקבות הפירוק, יש קבוצה לא קטנה של קצינים בכירים שנוטשים את צה''ל. בסופו של דבר תהליך הגיבוש של צה''ל הושלם, ויש כאן רמתכ''ל, אלופים שכפופים לדרג המדיני, וכו'. 
	
	אם חלק א' / תקופות 1-2 היו מלחמה פנימית, בה אין צבא סדיר מול צבא סדיר, בחלק ב' מדובר על מלחמה מול צבאות סדירים של מדינה שיש לה צבא סדיר (אבל עייף). המטרה של הצבאות הערבים הייתה לכבוש את המדינה מכל הצדדים (שתי גדות הירדן, מהנהר ועד הים). בחלק הזה מצליחים לדחוק את צה''ל ולהכנס פנימה, לכבוש שטחים. אחד מהמקומות אליהם הם פולשים וכובשים הוא קיבוץ ניצנים בסופו חברי ניצנים נכנעו ונלקחו לשבי. יחסי הכוחות, הם בצורה חד־משמעית נגדנו. דיברנו על זה, העדיפות העיקרית של צבאות ערב הייתה בציוד – מספר החיילים לא היה שונה בסדרי גודל. מה שהיה שונה היה הציוד הצבאי, בעיקר הכבד. יש להם מקלעים, תותחים, טנקים, חיל אוויר ומטוסים, וכו'. לנו יש טנק אחד בלי תותח. באמצעות העובדה שהמדינה הוקמה, היה אפשר לבצע גיוס חובה, להלוות מהאזרחים כסף למלחמה, להלאים רכבים, וכו'. 
	
	בהבט הדמוגרפי, עלו לכאן אנשים רבים. שליש מהכוח הלוחם בצה''ל בחלק השני של המלחמה היו עולים חדשים, חלקם אומנו בחו''ל ע''י המוסד לעליה ב'. עיקר הנשק הגיע מעסקת המשק הצ'כוסלובקית. צ'כוסלובקיה לא הייתה מוסרת מסמר לכאן ללא חסות ברה''מ. בגלל האינטרס של ברית המועצות, נמסר לכאן נשק צ'כוסלובקי. לאט לאט הגיע חיל האוויר, הגיעו חיילים, ונשקים כבדים. 
	
	צבאות מדינות ערב לא מיומנים במיוחד. חלק גדול מצה''ל נלחם במלחמת העולם השנייה – יש שם פרטיזנים, מורדים, רבים שנלחמו בצבאות בעלות הברית, מי שהיה ביישוב, אלפי אנשים בבריגדה, וכיוצא בזאת. צבאות מדינות ערב לא באמת נלחמו. 
	
	המיתוס הזה, של מעטים מול רבים, לא מתאים לכל המלחמה. גם הקטע של מלחמת קיום הוא לא אבסולוטי, ''מבצע עובדה הוא לא מלחמת קיום``. אבל אל תכתבו את זה במבחן. 
	
	בחודש הראשון, הוא תקופה 3 בחלק ב', יש את קרב ניצנים וכו'. בתומה יש הפוגה של חודש בלחימה עליה מחליט האו''ם, בה מדינת ישראל משקיעה מאמצים רבים בלהציב יעדים, להכין תוכניות לחימה, להפיל ספינות העונות לשם אלטלנה וכו', ומדינות ערב נחות. אחרי ההפוגה התרחש מבצע דני (במקור, מבצע לרלר), במסגרתו נכבשות רמלה ולוד, ושם יש גירוש במוצהר. במקור – לוד, רמלה, לתרון, רמלה (לרלר), אבל בסוף כבשו רק לר, בלי לר. המבצע אפשר מעבר לדרום ולירושלים. 
	
	אחרי הפוגה שנייה בתום לרלר התרחש מבצע עשרת הימים. במהלך ההפוגה בה היה אמור להתרחש מו''מ, תוך הפרתה, מדינת ישראל מבצעת 4 מבצעים במטרה להגיע לשולחן המו''מ במצב יותר טוב. במסגרת אחד המבצעים כובשים 14 ישובים בלבנון. ''מבצע עובדה הוא המבצע הכי מסובך ביקום, אף יריה לא נורתה שם, זה פשוט דוך בשני צירים בשני חיילות עד אילת, והמטרה הייתה להגיע לאילת ושהאיזור יהיה שלנו במסגרת המו''מ מול הירדנים``. גם ניתנה הנחייה שלא ירו ירייה אחת ויפגעו, כי בה בעת התרחש מו''מ ברודוס. 
	
\end{document}